\documentclass[a4paper,10.5pt]{article}

\usepackage[T1]{fontenc}
\usepackage[utf8]{inputenc}
\usepackage[brazil]{babel}
\usepackage[scaled=0.92]{helvet}
\renewcommand{\familydefault}{\sfdefault}
\usepackage[left=1cm,right=1cm,top=0.8cm,bottom=0.6cm]{geometry}
\usepackage{enumitem}
\usepackage{titlesec}
\usepackage{xcolor}
\usepackage{tabularx}
\usepackage[hidelinks]{hyperref}

%--------------------------------------------------------------
% Cores e estilos
%--------------------------------------------------------------
\definecolor{primary}{HTML}{1F3A5F}
\definecolor{muted}{HTML}{444444}

\hypersetup{colorlinks=true, urlcolor=primary}

\titleformat{\section}
  {\normalfont\large\bfseries\color{primary}}
  {}{0em}{}[\vspace{-0.7em}\color{primary}\rule{\textwidth}{0.8pt}]
\titlespacing*{\section}{0pt}{1.0em}{0.55em}

\setlist[itemize]{leftmargin=1.1em, topsep=2pt, itemsep=1.5pt, parsep=0pt}
\pagestyle{empty}
\setlength{\parindent}{0pt}

% Cabeçalho de experiência: {cargo}{empresa}{período}
\newcommand{\cargo}[3]{%
  \vspace{2pt}
  \begin{tabularx}{\textwidth}{@{}Xr@{}}
    \textbf{#1} \textnormal{\textbar{} #2} & \textcolor{muted}{#3}
  \end{tabularx}%
  \vspace{1pt}
}

% Linha de competência: {rótulo}{conteúdo}
\newcommand{\skillline}[2]{\textbf{#1:} #2\\}

\begin{document}

%==============================================================
% CABEÇALHO
%==============================================================
\begin{center}
  {\LARGE\bfseries\color{primary} ROBSON THIAGO DUARTE DA SILVA}\\[3pt]
  {\large Engenheiro de Front-End \textbar{} React \textbar{} TypeScript \textbar{} JavaScript}\\[5pt]
  \small
  Recife -- PE (disponível para trabalho presencial) \textbar{}
  \href{mailto:robgol.fidel@gmail.com}{robgol.fidel@gmail.com} \textbar{}
  \href{https://sites.google.com/cin.ufpe.br/rtds/p\%C3\%A1gina-inicial}{Portfólio}
\end{center}

\vspace{-0.3em}

%==============================================================
% RESUMO
%==============================================================
\section{Resumo Profissional}
Desenvolvedor Front-End com atuação profissional na construção de experiências digitais em produção para o
\textbf{AVAMEC}, plataforma educacional do Ministério da Educação utilizada por milhares de professores em todo o
Brasil. Trabalho diário com \textbf{JavaScript, TypeScript e React}, arquitetando \textbf{componentes reutilizáveis},
consumindo APIs REST e implementando fielmente mockups entregues pelo time de Design. Rotina consolidada de
\textbf{Git e fluxos colaborativos} (pull requests, rebase e code review), ambiente \textbf{Unix/linha de comando} e
atenção a \textbf{performance, acessibilidade e semântica} em HTML5/CSS3. Estudante de Engenharia da Computação na
UFPE, com forte histórico de entrega sob pressão --- \textbf{6 premiações em hackathons} e um produto homenageado no
\textbf{Prêmio Innovare (CNJ/STF)}. Perfil proativo, colaborativo e habituado a times multidisciplinares.

%==============================================================
% COMPETÊNCIAS
%==============================================================
\section{Competências Técnicas}
\skillline{Linguagens \& Tipagem}{JavaScript (ES6+), TypeScript, HTML5 semântico, CSS3 (Flexbox, Grid, responsividade), Python, SQL}
\skillline{Front-End}{React (hooks, composição de componentes, gerenciamento de estado), arquitetura de componentes reutilizáveis, design systems, consumo e integração de APIs REST, Web Components, noções de PWA}
\skillline{Fundamentos}{Estruturas de dados e algoritmos, padrões modernos da linguagem (módulos, assincronismo, imutabilidade), boas práticas de código limpo}
\skillline{Ecossistema \& Ferramentas}{Node.js/npm, bundlers e ferramentas de build, ESLint/Prettier, DevTools (debugger, auditorias de performance e acessibilidade via Lighthouse)}
\skillline{Versionamento \& Colaboração}{Git/GitHub, pull requests, rebase, code review entre pares, GitHub Actions (CI/CD), metodologias ágeis}
\skillline{Ambiente}{Linux/Unix e linha de comando, noções de Docker e desenvolvimento conteinerizado}
\skillline{Qualidade}{Acessibilidade (WCAG, navegação por teclado, contraste), performance web, noções de segurança em aplicações Web}
\skillline{Interesses técnicos}{Open source, Internet das Coisas (IoT) com ESP32, redes intraveiculares, inovação aplicada ao cotidiano}

%==============================================================
% EXPERIÊNCIA
%==============================================================
\section{Experiência Profissional}

\cargo{Desenvolvedor Front-End}{V-Lab UFPE --- Projeto AVAMEC/MEC}{Atual --- 1 ano e 5 meses}
\begin{itemize}
  \item Desenvolvimento e manutenção de experiências digitais em produção para o \textbf{AVAMEC}, ambiente virtual de aprendizagem do Ministério da Educação, utilizando \textbf{React e TypeScript}.
  \item Arquitetura de \textbf{componentes de front-end reutilizáveis}, aproveitados por outros integrantes do time e padronizando comportamento e estilo entre telas.
  \item Implementação fiel de \textbf{mockups do time de Design}, com alinhamento constante sobre requisitos, viabilidade técnica e alternativas de solução.
  \item Integração de interfaces com \textbf{APIs REST}, tratando estados de carregamento, erro e consistência dos dados entre front-end e serviços de back-end.
  \item Participação ativa em \textbf{code review}, sugerindo melhorias e incorporando feedback dos colegas; versionamento em Git com fluxo de branches e pull requests.
  \item Colaboração com equipe multidisciplinar (Design, gestão de projeto e back-end) em ciclos ágeis de entrega.
\end{itemize}

\cargo{Tutor de Tecnologia}{CESAR School --- Projeto Florescendo Talentos}{Atual --- 2 anos e 3 meses}
\begin{itemize}
  \item Capacitação de estudantes de Ensino Médio e EJA em \textbf{desenvolvimento web (HTML, CSS e JavaScript)}, \textbf{lógica de programação com Python} e manipulação de dados com SQL e Google Sheets.
  \item Construção de materiais didáticos e projetos práticos, exigindo domínio conceitual dos fundamentos de linguagem e capacidade de explicar decisões técnicas com clareza --- competência aplicada diretamente em code reviews e discussões de requisitos.
  \item Mentoria individual e acompanhamento de evolução, reforçando cultura de aprendizado compartilhado.
\end{itemize}

\cargo{Analista de Help Desk Júnior}{Datamétrica}{1 ano}
\begin{itemize}
  \item Suporte de TI Nível 1 a órgãos do Governo do Estado de Pernambuco, aplicando o framework \textbf{ITIL} na gestão de incidentes e requisições.
  \item Diagnóstico e depuração de problemas em ambientes heterogêneos, incluindo sistemas operacionais \textbf{Unix/Linux} e uso de linha de comando.
  \item Comunicação técnica com usuários finais e times internos, traduzindo problemas de negócio em descrições técnicas acionáveis.
\end{itemize}

\cargo{Professor de Informática Básica}{News Center Cursos \textbar{} ONG Coque Connecta (voluntário)}{10 meses}
\begin{itemize}
  \item Ensino de informática e ferramentas digitais para turmas regulares e para crianças da comunidade do Coque (Recife), com foco em inclusão digital.
\end{itemize}

\cargo{Jovem Aprendiz --- Administrativo}{Proxxi Tecnologia}{1 ano e 6 meses}
\begin{itemize}
  \item Organização e controle de documentos com adequação aos requisitos da \textbf{LGPD} e suporte a setores internos.
\end{itemize}

%==============================================================
% PROJETOS
%==============================================================
\section{Projetos e Reconhecimentos}

\cargo{Desenvolvedor Full Stack --- Fênix Connect}{Homenageado no Prêmio Innovare (21ª ed.)}{2024}
\begin{itemize}
  \item Aplicativo voltado à reintegração social de egressos do sistema prisional, \textbf{homenageado no Prêmio Innovare} --- iniciativa do Instituto Innovare em parceria com \textbf{CNJ, STF e Grupo Globo}.
  \item Responsável pela interface do usuário e pela integração com o back-end: construção de componentes, modelagem de dados, consumo de API e testes de validação com usuários reais.
  \item Produto concebido com foco em usabilidade e acessibilidade para um público de baixa familiaridade digital.
\end{itemize}

\textbf{Premiações em Hackathons} --- 6 colocações no pódio, com entrega de protótipos funcionais em prazos curtos:
\begin{itemize}
  \item \textbf{1º lugar} --- Hacker Cidadão 9.0 \textbar{} \textbf{1º lugar} --- 1º Hackathon de SUAPE (2024)
  \item \textbf{2º lugar} --- Hacker Cidadão 7.0 \textbar{} \textbf{2º lugar} --- Ideatel 2025 (UFRPE)
  \item \textbf{2º lugar} --- Hackathon BBTS 2025 \textbar{} \textbf{2º lugar} --- EDscript Hackathon do EDS (2026)
\end{itemize}

%==============================================================
% FORMAÇÃO
%==============================================================
\section{Formação Acadêmica}
\cargo{Bacharelado em Engenharia da Computação}{Universidade Federal de Pernambuco (UFPE)}{Conclusão prevista: dez/2027}

%==============================================================
% IDIOMAS
%==============================================================
\section{Idiomas}
Português (nativo) \textbar{} Inglês técnico --- leitura e escrita de documentação técnica, discussão de funcionalidades e bugs

%==============================================================
% COMUNIDADE
%==============================================================
\section{Comunidade e Voluntariado}
\begin{itemize}
  \item \textbf{Eventos de tecnologia:} Rec'n'Play --- Arena IA e Arena de Inovação (2024 e 2025); NASA Space Apps (2025); Secomp UFPE (2024); Agile Trends NE (2025, diretor de palco); Casa do Julgamento (2025).
  \item \textbf{Ação social:} Anjos da Noite (2021--2022) e Programa Compaixão / ONG O Portal (2024--2025) --- distribuição de alimentos e roupas a pessoas em situação de rua no Recife.
\end{itemize}

\end{document}
